\scanpage{237}
\mspage{121}

\diagram{\begin{tikzcd}[ampersand replacement=\&]
{(54)}\&1\arrow[r]\&G_0\arrow[r]\&G\arrow[r,"\delta"]\&Z\\
\&\&\&Z\arrow[u,hook]\arrow[ur,"{\lambda\mapsto\lambda^2}"']\&{,}
\end{tikzcd}}

où $\delta$ (jouant le rôle d'un \textquotedblleft{}déterminant\textquotedblright) exprime l'opération de $G$ sur $\pi$ par translations, et où on a posé

\[
\begin{aligned}
G_0=\mathrm{Ker}\,\delta&=\{g\in\deleted{\mathcal G}\mathcal A\mid g(\varphi)=\varphi\quad\forall\varphi\in\pi\}\\
&=\left\{g\in\mathcal A\ \middle|\ \substack{\deleted{\text{i.e.}}\ g\in\mathcal G\ \text{i.e. }g\text{ centralise }Z\\g(\varphi_0)=\varphi_0}\right\}.
\end{aligned}
\tag*{(55)}
\]

\margin{\uline{NB} Si $R$ séparée en $\check S$, on a $(51')\ \left\{\begin{aligned}G_0&\xrightarrow{\sim}\text{groupe des automorphismes}\\[-2pt]&\text{du graphe }\gamma_{\varphi_0}=(S,A_{\varphi_0})\\ \lambda&\longmapsto\lambda_S\quad\text{qui commutent à }Z_S.\end{aligned}\right.$}

Le fait que $Z\subset G$, et la commutativité exprimée dans (54)

\[
\delta(\lambda)=\lambda^2\qquad\forall\lambda\in Z
\tag*{(56)}
\]

provient de la relation

\[
\lambda(\varphi)=\lambda_{\check S}\circ\underbrace{\varphi\circ\lambda_S^{-1}}_{\lambda_{\check S}\varphi}=\lambda_{\check S}^{2}\circ\varphi=T_{\lambda^2}(\varphi).
\]

La définition de $\delta(g)$ pour $g\in G$ s'explicite par

\[
g(\varphi)=T_{\delta(g)}\varphi\ \bigl(=\delta(g)_{\check S}\circ\varphi=\varphi\circ\delta(g)_S\bigr)\qquad\substack{\forall\varphi\in\pi\\g\in G}
\tag*{(57)}
\]

qui caractérise \deleted{\unsure{$\forall$}} $\delta(g)$, pour \uline{un} $\varphi=\varphi_0$ fixé déjà.

On a les Conditions équivalentes sur $M=(\Gamma,Z,\pi)$

\[
\begin{aligned}
&\text{a) L'hom }\delta\text{ est surjectif}\\
&\text{b) }\deleted{\unsure{\mathcal G}}G\text{ opère transitivement sur }\pi\\
&\text{c) }\forall\lambda\in Z,\ \deleted{\unsure{\lambda\varphi_0}}\ \text{la structure de}\\
&\quad\text{graphes définie par } A_\lambda\subset\underset{\{1,2\}}{\mathfrak P}(S)\\
&\quad A_\lambda=\{\{s,\lambda t\}\mid\{s,t\}\in S,\ \text{t.q. }\{s,\deleted{\unsure{\lambda}}t\}\inA_0\}\\
&\quad\text{est isomorphe à }\unsure{\gamma}(S,A_0),\ \text{par une}\\
&\quad\deleted{\unsure{f}}\ f\in\unsure{\mathfrak S|\mathcal G}_S\text{ qui centralise }Z_S\\
&\text{d) }\inserted{\forall\lambda\in Z}\ \exists f_\lambda\in\mathfrak S_S,\ \text{centralisant }Z_S,\ \text{telle que}\\
&\quad\forall s,t\in S,\ \text{on a }\{s,t\}\in A_0\Longleftrightarrow\{f_\lambda^{-1}s,f_\lambda^{-1}\lambda t\}\in A_0
\end{aligned}
\tag*{(58)}
\]

% slanted
\margin{ici on suppose choisi un $\varphi_0\in\pi$, d'où une structure de graphes --- et on suppose $R$ \deleted{\unsure{non}} séparée en $\check S$}

On dit que la condition d'homogénéité est satisfaite

% archive sheet number: 236
